\chapter*{Введение}
\addcontentsline{toc}{chapter}{Введение}
\label{ch:intro}

Актуальность темы обусловлена ростом числа территориально распределённых информационных систем и вместе с этим увеличение количества сетевых угроз.
В связи с этим информационные системы требуют применения современных средств сегментации, маршрутизации, мониторинга и защиты сетевого трафика.

Целью данной научно-исследовательской работы является практика сетевого инженеринга на виртуальной модели - траблшутинг, построение сетей и основ безопасности вместе с атакующими техниками и анализом трафика для дальнейшего поддержания безопасности в реальных сетевых инфраструктурах.

Основные задачи исследования, решаемые в рамках работы:

\begin{enumerate}
    \item TODO
    \item TODO
    \item TODO
\end{enumerate}

Объектом исследования являются локальные вычислительные сети. Предметом исследования выступает TODO.

Положения, выносимые на защиту:

1) TODO

2) TODO

3) TODO

Достоверность и обоснованность полученных результатов подтверждается по нескольким пунктам:

1) TODO

2) TODO

3) TODO


Теоретическая значимость работы заключается в TODO.

Практическая значимость состоит в TODO.

Планируемый научно-технический уровень разработки соответствует TODO.

\textbf{Объем и структура работы.} Отчёт состоит из введения, заключения, \total{chapter} глав, списка использованной литературы (0 источников). Содержит \ztotpages\ страниц текста, включая \totalfigures\ рисунка и \totaltables\ таблиц, 0 формул.

Первая глава посвящена TODO. Вторая - TODO. В третьей главе представлен TODO.
